\documentclass[11pt,a4paper]{article}
% Estilo comum dos PDFs individuais da Sprint 1.
\usepackage[utf8]{inputenc}
\usepackage[T1]{fontenc}
\usepackage[brazilian]{babel}
\usepackage[a4paper,margin=2.35cm]{geometry}
\usepackage{xcolor}
\usepackage{hyperref}
\usepackage{enumitem}
\usepackage{listings}
\usepackage{parskip}
\usepackage{needspace}
\definecolor{azul}{RGB}{20,63,112}
\hypersetup{colorlinks=true,linkcolor=azul,urlcolor=azul}
\setlist[itemize]{leftmargin=*,itemsep=3pt}
\lstset{basicstyle=\ttfamily\small,frame=single,breaklines=true,columns=fullflexible}
\sloppy
\newcommand{\codigo}[1]{\texttt{#1}}
\newcommand{\titulo}[2]{\begin{center}{\LARGE\bfseries #1\par}\vspace{4pt}{\large #2\par}\vspace{8pt}\rule{0.8\textwidth}{0.4pt}\par\vspace{5pt}{\small Matheus Moreli Durão Silva -- Matrícula 12621CBS217\par Cibersegurança - Bacharelado - Noturno\par}\end{center}}
\newcommand{\cardfonte}[1]{\par\smallskip\noindent\textit{Referência: }\href{#1}{Trello - enunciado e checklist do card}.}

\begin{document}\titulo{S1\_05 - Variáveis e tipos de dados}{Guardar informações de um ativo}
\section*{O que o card pede}
O entregável correto solicita ao usuário o ID, nome e responsável de um ativo de TI e mostra essas informações formatadas. Ele demonstra variáveis, entrada, saída e tipos de dados; não é apenas uma declaração abstrata de variáveis.
\section*{Como funciona}
Uma variável é um nome que referencia um valor enquanto o programa roda. \codigo{nome = "Firewall"} guarda texto, portanto \codigo{nome} é \codigo{str}. IDs e quantidades normalmente são \codigo{int}; preços podem ser \codigo{float}; uma resposta de sim/não pode ser \codigo{bool}. \codigo{input()} sempre devolve texto, então \codigo{int(input("ID: "))} faz a conversão explícita para número inteiro.
\begin{lstlisting}
id_ativo = int(input("ID do ativo: "))
nome = input("Nome: ")
responsavel = input("Responsavel: ")
print(f"ID: {id_ativo} | Ativo: {nome} | Responsavel: {responsavel}")
\end{lstlisting}
\section*{Aplicação no projeto}
No \codigo{app.py}, cada cadastro começa com valores recebidos do usuário e guardados em variáveis. Depois, eles formam um dicionário que será salvo em JSON. Dar nomes claros, como \codigo{responsavel} em vez de \codigo{x}, reduz erros e permite explicar a lógica.
\section*{Checklist demonstrado}
Conceito de variável, tipos primários, conversões, entrada/saída e aplicação de dados de ativos foram praticados. Para testar, informar um ID numérico e conferir se os três dados aparecem na mesma mensagem. Este é o entregável do card; exemplos adicionais só servem para estudar, não o substituem.
\cardfonte{https://trello.com/c/DmvtABXt}
\end{document}
